\appendixsection[KW]{Auswertung Kreislaufwirtschaft (Raoul Bunschoten)}
\label{anlage:interview-bunschoten}

Dieser Anhang fasst zentrale Erkenntnisse aus dem Gespräch mit Raoul Bunschoten zusammen. Im Mittelpunkt stehen dabei nicht technische Einzelfunktionen, sondern übergeordnete Fragen: Wie können fragmentierte Prozesse verbunden werden? Welche Rolle spielen unterschiedliche Akteure? Wie lassen sich unvollständige Daten, Szenarien und dynamische Planungssituationen in einem digitalen Werkzeug abbilden?

\subsection{Fehlende Kontinuität als zentrales Problem}

Ein wiederkehrendes Thema im Gespräch ist die fehlende Kontinuität entlang der Wertschöpfungskette. Bunschoten beschreibt, dass viele Prozesse zwischen Materialherkunft, Planung, Produktion, Regulierung, Nutzung, Rückbau und Wiederverwendung nicht ausreichend miteinander verbunden sind.

\begin{Blockquote}
There are many gaps in the continuity of the workflow.
\end{Blockquote}

Für die Entwicklung eines Reuse-Werkzeugs bedeutet das: Die Herausforderung liegt nicht nur darin, Materialdaten zu sammeln. Entscheidend ist, wie Informationen zwischen unterschiedlichen Prozessschritten weitergegeben, übersetzt und nutzbar gemacht werden können.

Ein digitales Werkzeug sollte daher nicht nur als Materialkatalog funktionieren, sondern als verbindende Struktur zwischen Akteuren, Datenformaten und Entscheidungsphasen.

\subsection{Ein Tool als Teil eines größeren Ökosystems}

Bunschoten beschreibt das Value Chain Digital Model nicht als einzelnes abgeschlossenes Werkzeug, sondern als Plattform beziehungsweise als Zusammenspiel mehrerer Werkzeuge.

\begin{Blockquote}
It's only a family of tools, right? It's not one tool.
\end{Blockquote}

Diese Perspektive ist für ein digitales Reuse-Werkzeug besonders wichtig. Wiederverwendung betrifft viele Ebenen gleichzeitig: Materialinventar, Bewertung, Simulation, Entwurf, Fertigung, Logistik, Regulierung, Beschaffung und spätere Wiederverwendung.

Das Werkzeug soll als Teil eines offenen Ökosystems funktionieren. Es verbindet unterschiedliche Datenquellen, Bewertungsverfahren und Akteursgruppen.

Die zentrale Qualität eines solchen Systems liegt nicht darin, alles selbst zu lösen, sondern relevante Schnittstellen herzustellen.

\subsection{Wiederverwendung als Verhandlungsprozess}

Ein weiterer wichtiger Aspekt ist, dass Wiederverwendung nicht als rein technisches Problem verstanden werden kann. Sie ist ein Aushandlungsprozess zwischen vielen Beteiligten: Materiallieferant:innen, Planer:innen, Ingenieur:innen, Fertiger:innen, Bauherr:innen, Behörden, Rückbauunternehmen, Nutzer:innen und politischen Entscheidungsträger:innen.

Diese Akteure arbeiten nicht nur mit unterschiedlichen Daten, sondern auch mit unterschiedlichen Begriffen, Prioritäten und Bewertungsmaßstäben. Was für eine Planerin eine gestalterische Chance ist, kann für einen Ingenieur ein Risiko, für eine Bauherrin eine Kostenfrage und für eine Behörde ein Nachweisthema sein.

Ein Reuse-Werkzeug sollte deshalb nicht nur Informationen darstellen, sondern gemeinsame Bewertung und Kommunikation ermöglichen. Es sollte eine Grundlage schaffen, auf der verschiedene Akteure dieselbe Situation diskutieren und Entscheidungen nachvollziehbar aushandeln können.

\subsection{Szenarien statt automatischer Lösungen}

Bunschoten betont, dass das VCDM nicht als Werkzeug verstanden werden sollte, das fertige Lösungen erzeugt. Vielmehr dient es dazu, Szenarien sichtbar und vergleichbar zu machen.

\begin{Blockquote}
It doesn't help you solve things, you also have to make your own decisions as stakeholders.
\end{Blockquote}

Für die Tool-Entwicklung ist dieser Punkt zentral. Ein Werkzeug für wiederverwendete Materialien sollte nicht versuchen, automatisch die „richtige“ Lösung zu erzeugen. Reuse-Entscheidungen hängen von vielen Faktoren ab: Materialverfügbarkeit, Zustand, Kosten, CO\textsubscript{2}-Wirkung, technische Eignung, regulatorische Anforderungen, Logistik, gestalterische Qualität und zukünftige Wiederverwendbarkeit.

Das Werkzeug sollte deshalb Entscheidungsräume öffnen. Es sollte Alternativen vergleichbar machen und sichtbar zeigen, auf welchen Annahmen, Daten und Unsicherheiten eine bestimmte Option basiert.

\subsection{Arbeiten mit unvollständigem Wissen}

Bei wiederverwendeten Materialien sind Daten häufig unvollständig. Herkunft, frühere Belastung, Zustand, Qualität, Verfügbarkeit oder zukünftige Nutzbarkeit eines Bauteils sind oft nur teilweise bekannt. Bunschoten weist darauf hin, dass digitale Daten allein nicht ausreichen, um komplexe reale Situationen vollständig zu erfassen.

\begin{Blockquote}
The pure data don't tell you full story.
\end{Blockquote}

Für ein Reuse-Werkzeug bedeutet das, dass Unsicherheit nicht versteckt werden sollte. Stattdessen muss sichtbar werden, welche Informationen gesichert sind, welche auf Annahmen beruhen und wo zusätzliche Prüfung oder Expert:innenwissen notwendig ist.

Ein gutes Modell sollte also nicht nur Daten speichern, sondern unterschiedliche Wissenszustände abbilden können. Gerade diese Fähigkeit unterscheidet ein entscheidungsunterstützendes Werkzeug von einem einfachen digitalen Inventar.

\subsection{Planung als dynamischer Prozess}

Bunschoten beschreibt Planung nicht als statisches Abbild, sondern als dynamischen Prozess. Projektbedingungen verändern sich durch neue Vorschriften, politische Entscheidungen, Materialverfügbarkeiten, Kostenentwicklungen, Zeitdruck oder wechselnde Akteurskonstellationen.

\begin{Blockquote}
Not a still image of something, but have a dynamic image.
\end{Blockquote}

Für wiederverwendete Materialien ist diese Dynamik besonders relevant. Ein Bauteil kann verfügbar sein und später nicht mehr zugänglich sein. Eine Prüfung kann neue Ergebnisse liefern. Eine Vorschrift kann sich ändern. Ein Projektziel kann neu bewertet werden.

Ein Reuse-Werkzeug sollte deshalb nicht nur einen festen Zustand dokumentieren, sondern Veränderung als Teil des Planungsprozesses berücksichtigen. Es muss aktualisierbar bleiben und auf neue Informationen reagieren können.

\subsection{Anpassungsfähigkeit als zentrale Anforderung}

Aus der Dynamik von Planung folgt, dass ein digitales Reuse-Werkzeug flexibel auf veränderliche Bedingungen reagieren können muss. Ein Materialinventar darf nicht als statische Grundlage verstanden werden. Es muss erweitert, korrigiert, aktualisiert und mit neuen Informationen verknüpft werden können.

Das Werkzeug sollte zum Beispiel sichtbar machen können, was passiert, wenn ein Material nicht mehr verfügbar ist, ein Bauteil zusätzliche Prüfung benötigt, eine Vorschrift angepasst wird oder eine andere Wiederverwendungsstrategie sinnvoller erscheint.

Das Werkzeug unterstützt damit sowohl den Entwurf als auch den Umgang mit veränderten Projektbedingungen.
